\section{Final Problem Statement}
\label{sec:final-problem-statement}
The preceding review shows that adaptive reading is an active research area with a clear unmet need.
Prior work in the Reading the Reader project produced capable but tightly coupled prototypes \cite{refsgaard2024rtr, pereira2024typography, palinec2024reactive}, the literature establishes that intervention design measurably affects reading recovery and preference \cite{jensen2025context}, and the market survey confirms that no existing product unites real-time gaze sensing, automatic typographic adaptation, and researcher-facing experiment control in a single modular platform.
The gap analysis in Sec.~\ref{sec:gaps-opportunities} reduced this to three concerns: experimentability, the coupling of the sensing layer, and the integration of external decision modules.

This thesis therefore addresses the following problem.
\begin{quote}
How can a researcher-operated adaptive reading platform be architected so that sensing, eye-movement analysis, decision strategy, and intervention execution are independently replaceable, while supporting real Tobii-backed reading sessions and applying context-preserving micro-interventions without disrupting the participant's reading flow?
\end{quote}

We decompose this into four research questions, each of which maps to evaluation criteria developed in \cref{ch:evaluation}.

\begin{description}
    \item[RQ1 (Modularity).] How can the platform separate sensing, analysis, decision, and intervention into modules with stable contracts, such that a new intervention or decision provider can be added without modifying the core reading runtime?
    \item[RQ2 (Sensing pipeline).] How can a real Tobii eye-tracking stream be converted into oculomotor events (fixations, saccades, and regressions) within a latency budget that preserves live adaptation during a reading session, in a manner that is inspectable and reproducible? The associated latency and feasibility criteria are defined and measured in \cref{ch:evaluation}.
    \item[RQ3 (Context-preserving intervention).] How can typographic micro-interventions be committed at controlled boundaries so that they adapt the text while preserving the reader's context, consistent with the recovery concerns identified in the literature \cite{jensen2025context}?
    \item[RQ4 (Researcher control and experimentability).] How can the platform give the researcher the control and the auditable record (including session replay) needed to operate experiments and compare interventions and decision strategies?
\end{description}

These questions deliberately emphasise architecture and runtime feasibility over a controlled human-subjects effect study, in line with the thesis scope: the contribution is a defendable, modular platform on which interventions and external, AI-driven decision providers can be compared in future work, rather than an end-to-end empirical result for any single intervention.
